\subsection{べき関数}
まずは $x^n$ を定義から微分しよう.正の整数 $n\ge2$ について二項定理を使うと,
\[
 \frac{(x+h)^n-x^n}{h}
 =nx^{n-1}+\sum_{k=2}^n\binom nk x^{n-k}h^{k-1}
\]
となる.ここで $0\le k\le n$ の整数 $k$ について,\,$\binom nk=\frac{n!}{k!(n-k)!}$ は二項係数であり,\,$n!=1\cdot2\cdots n$,\,$0!=1$ とする.右辺の和の各項は $h$ を因数にもつので,\,$h\to0$ とすれば
\[
 \frac{\dd}{\dd x}x^n=nx^{n-1}
\]
を得る.\,$n=1$ の場合は差商が $1$ なので,この式は $n=1$ でも成り立つ.定数の微分は $0$ である.

負の整数べきは,逆数の微分から同じ形になる.\,$m$ を正の整数とし,\,$x\ne0$ とすると,
\begin{align*}
 (x^{-m})'&=\left(\frac1{x^m}\right)'\\
 &=-\frac{mx^{m-1}}{x^{2m}}\\
 &=-mx^{-m-1}
\end{align*}
となり,\,$n=-m$ と置けば $nx^{n-1}$ の形である.実数べきへの拡張は,対数関数を微分した後で説明する.
